% !TEX program = pdflatex
% ==========================================================
% PLANTILLA BEAMER PROFESIONAL
% Universidad Santo Tomás
% Luis Tulio González Alcaino
% ==========================================================

\newif\ifhandout
%\handoutfalse
\handouttrue

\ifhandout
\documentclass[aspectratio=169,10pt,handout]{beamer}
\else
\documentclass[aspectratio=169,10pt]{beamer}
\fi

% ----------------------------------------------------------
% PAQUETES
% ----------------------------------------------------------

\usepackage[spanish,es-noshorthands]{babel}
\usepackage[T1]{fontenc}
\usepackage[utf8]{inputenc}

\usepackage{lmodern}
\usepackage{amsmath,amsfonts,amssymb}
\usepackage{ragged2e}
\usepackage{enumitem}
\usepackage{multicol}

\usepackage{tikz}
\usetikzlibrary{positioning}

\usepackage{fontawesome5}

% ----------------------------------------------------------
% COLORES INSTITUCIONALES
% ----------------------------------------------------------

\definecolor{USTBlue}{RGB}{0,70,127}
\definecolor{USTBlueDark}{RGB}{0,45,85}
\definecolor{USTTeal}{RGB}{0,153,117}
\definecolor{USTGray}{RGB}{120,120,120}
\definecolor{USTSoft}{RGB}{248,249,250}
\definecolor{USTText}{RGB}{35,40,45}
\definecolor{WarmGray}{RGB}{120,120,120}

% ----------------------------------------------------------
% TEMA BASE
% ----------------------------------------------------------

\usetheme{Madrid}
\usecolortheme[named=USTBlue]{structure}
\setbeamertemplate{navigation symbols}{}

% ----------------------------------------------------------
% COLORES DE PIE
% ----------------------------------------------------------

\setbeamercolor{author in head/foot}{bg=USTSoft,fg=USTGray}
\setbeamercolor{date in head/foot}{bg=USTSoft,fg=USTGray}
\setbeamercolor{title in head/foot}{bg=USTSoft,fg=USTGray}

% ----------------------------------------------------------
% PIE DE PÁGINA
% ----------------------------------------------------------

\setbeamertemplate{footline}{%
	\leavevmode%
	\hbox{%
		\begin{beamercolorbox}[wd=.85\paperwidth,ht=2.5ex,dp=1ex,leftskip=0.4cm]{author in head/foot}%
			\scriptsize
			\textbf{Fundamentos de Matemática Básica}
			\quad|\quad
			Agronomía
			\quad|\quad
			Universidad Santo Tomás Talca
		\end{beamercolorbox}%
		\begin{beamercolorbox}[wd=.15\paperwidth,ht=2.5ex,dp=1ex,rightskip=0.3cm plus1fil]{date in head/foot}%
			\scriptsize \insertframenumber/\inserttotalframenumber
		\end{beamercolorbox}%
	}%
}

% ----------------------------------------------------------
% BLOQUES PEDAGÓGICOS
% ----------------------------------------------------------

\setbeamercolor{block title}{bg=USTBlue!15,fg=USTBlueDark}
\setbeamercolor{block body}{bg=USTSoft,fg=black}

\setbeamercolor{block title example}{bg=USTTeal!20,fg=USTBlueDark}
\setbeamercolor{block body example}{bg=USTSoft,fg=black}

\setbeamercolor{block title alerted}{bg=red!15,fg=black}
\setbeamercolor{block body alerted}{bg=red!3,fg=black}

% ----------------------------------------------------------
% COMANDOS Y ENTORNOS AUXILIARES
% ----------------------------------------------------------

\newcommand{\destacar}[1]{\textcolor{USTBlueDark}{\textbf{#1}}}
\newcommand{\alerta}[1]{\textcolor{red!70!black}{\textbf{#1}}}
\newcommand{\pp}{\pause}

\newenvironment{metaBox}
{\begin{block}{Meta de aprendizaje}}
	{\end{block}}

\newenvironment{ideaBox}
{\begin{exampleblock}{Idea clave}}
	{\end{exampleblock}}

% ----------------------------------------------------------
% PORTADA PROFESIONAL PREMIUM
% ----------------------------------------------------------
\newcommand{\PortadaProfesionalUST}{%
	\begin{frame}[plain]
		\begin{tikzpicture}[remember picture,overlay]
			
			\fill[white] (current page.south west) rectangle (current page.north east);
			\fill[USTBlueDark] (current page.north west) rectangle ([yshift=-1.35cm]current page.north east);
			\fill[USTTeal] ([yshift=-1.35cm]current page.north west) rectangle ([yshift=-1.52cm]current page.north east);
			\fill[USTSoft] (current page.south west) rectangle ([xshift=2.15cm]current page.north west);
			
			\fill[USTBlue!10] ([xshift=1.05cm,yshift=-1.8cm]current page.north west) circle (0.72cm);
			\fill[USTTeal!12] ([xshift=1.65cm,yshift=-4.3cm]current page.north west) circle (0.42cm);
			\fill[USTBlue!6] ([xshift=-1.15cm,yshift=1.0cm]current page.south east) circle (1.15cm);
			
			\fill[black!6,rounded corners=7pt]
			([xshift=-4.95cm,yshift=-3.25cm]current page.center) rectangle
			([xshift=4.95cm,yshift=2.45cm]current page.center);
			
			\fill[white,rounded corners=7pt]
			([xshift=-5.1cm,yshift=-3.1cm]current page.center) rectangle
			([xshift=4.8cm,yshift=2.6cm]current page.center);
			
			\node[anchor=north east,text=white]
			at ([xshift=-0.65cm,yshift=-0.48cm]current page.north east)
			{\small\bfseries Universidad Santo Tomás\quad Ciencias Básicas Talca};
			
			\fill[USTTeal] ([xshift=-5.1cm,yshift=-3.1cm]current page.center) rectangle
			([xshift=-4.82cm,yshift=2.6cm]current page.center);
			
			\node[align=center] at ([xshift=0.15cm,yshift=-0.1cm]current page.center){%
				\begin{minipage}{0.66\paperwidth}
					\centering
					{\Large\bfseries\color{USTText}\insertauthor\par}
					\vspace{0.95cm}
					{\fontsize{22}{26}\selectfont\bfseries\color{USTBlueDark}\inserttitle\par}
					\vspace{0.28cm}
					{\large\color{USTTeal}\insertsubtitle\par}
					\vspace{0.95cm}
					{\small\color{USTGray}\insertdate\par}
				\end{minipage}
			};
			
			\draw[USTGray!35,line width=0.35pt]
			([xshift=0.95cm,yshift=0.72cm]current page.south west) --
			([xshift=-0.95cm,yshift=0.72cm]current page.south east);
			
			\node[anchor=south,text=USTGray]
			at ([yshift=0.2cm]current page.south)
			{\scriptsize
				\textbf{Fundamentos de Matemática Básica}
				\;|\;
				Agronomía
				\;|\;
				Universidad Santo Tomás Talca};
			
		\end{tikzpicture}
	\end{frame}
}

% ----------------------------------------------------------
% DATOS DEL CURSO
% ----------------------------------------------------------

\title{Fundamentos de Matemática Básica}
\subtitle{\Large{Clase: Potencias}}
\author{Prof. Luis González Alcaino}
\institute{Universidad Santo Tomás}
\date{Año académico 2026}

% ==========================================================
% DOCUMENTO
% ==========================================================

\begin{document}
	
	\PortadaProfesionalUST
	
	% ----------------------------------------------------------
	\begin{frame}{Contenidos de la clase}
		\begin{block}{Aprendizajes de la sesión}
			\begin{enumerate}
				\item Comprender el concepto de \destacar{potencia}.
				\item Interpretar la \destacar{base} y el \destacar{exponente}.
				\item Aplicar las \destacar{propiedades de las potencias}.
				\item Reconocer \destacar{patrones numéricos}.
				\item Resolver \destacar{problemas contextualizados en Agronomía}.
			\end{enumerate}
		\end{block}
	\end{frame}
	
	% ----------------------------------------------------------
	\begin{frame}{Concepto de potencia}
		\begin{block}{Definición}
			Una \destacar{potencia} representa una multiplicación repetida de un mismo factor:
			\[
			a^n = \underbrace{a\cdot a\cdot a\cdots a}_{n \text{ veces}}
			\]
			donde:
			\begin{itemize}
				\item \(a\) es la \destacar{base},
				\item \(n\) es el \destacar{exponente}.
			\end{itemize}
		\end{block}
		
		\pp
		
		\begin{exampleblock}{Ejemplos}
			\[
			2^4 = 2\cdot2\cdot2\cdot2 = 16
			\qquad
			3^3 = 27
			\]
			\[
			10^3 = 1000
			\qquad
			5^2 = 25
			\]
		\end{exampleblock}
	\end{frame}
	
	% ----------------------------------------------------------
	\begin{frame}{¿Por qué estudiar potencias en Agronomía?}
		\begin{columns}[T,totalwidth=\textwidth]
			\begin{column}{0.52\textwidth}
				\begin{block}{Situaciones reales}
					\begin{itemize}
						\item Número de semillas por hectárea.
						\item Crecimiento de bacterias o hongos.
						\item Concentración de nutrientes.
						\item Escalas muy grandes o muy pequeñas.
					\end{itemize}
				\end{block}
			\end{column}
			
			\begin{column}{0.44\textwidth}
				\begin{exampleblock}{Ejemplo rápido}
					\[
					0{,}00045\ \text{kg}=4{,}5\times10^{-4}\ \text{kg}
					\]
					La forma con potencias de 10 facilita cálculos y comparaciones.
				\end{exampleblock}
			\end{column}
		\end{columns}
		
		\vspace{0.2cm}
		\centering
		{\small \textcolor{WarmGray}{Las potencias ayudan a describir crecimiento, disminución y escalas técnicas.}}
	\end{frame}
	
	% ----------------------------------------------------------
	\begin{frame}{Propiedades de las potencias I}
		\begin{block}{1. Producto de potencias de igual base}
			\[
			a^m\cdot a^n = a^{m+n}
			\]
		\end{block}
		
		\begin{exampleblock}{Ejemplo}
			\[
			2^3\cdot2^4 = 2^{3+4}=2^7=128
			\]
		\end{exampleblock}
		
		\pp
		
		\begin{block}{2. Cociente de potencias de igual base}
			\[
			\frac{a^m}{a^n}=a^{m-n}, \qquad a\neq 0
			\]
		\end{block}
		
		\begin{exampleblock}{Ejemplo}
			\[
			\frac{5^6}{5^2}=5^{6-2}=5^4=625
			\]
		\end{exampleblock}
	\end{frame}
	
	% ----------------------------------------------------------
	\begin{frame}{Propiedades de las potencias II}
		\begin{block}{3. Potencia de una potencia}
			\[
			(a^m)^n = a^{m\cdot n}
			\]
		\end{block}
		
		\begin{exampleblock}{Ejemplo}
			\[
			(3^2)^4 = 3^{2\cdot4}=3^8=6561
			\]
		\end{exampleblock}
		
		\pp
		
		\begin{block}{4. Potencia de un producto}
			\[
			(ab)^n = a^n b^n
			\]
		\end{block}
		
		\begin{exampleblock}{Ejemplo}
			\[
			(2\cdot5)^3 = 2^3\cdot5^3 = 8\cdot125 = 1000
			\]
		\end{exampleblock}
	\end{frame}
	
	% ----------------------------------------------------------
	\begin{frame}{Propiedades de las potencias III}
		\begin{block}{5. Potencia de un cociente}
			\[
			\left(\frac{a}{b}\right)^n = \frac{a^n}{b^n}, \qquad b\neq 0
			\]
		\end{block}
		
		\begin{exampleblock}{Ejemplo}
			\[
			\left(\frac{2}{3}\right)^3 = \frac{2^3}{3^3}=\frac{8}{27}
			\]
		\end{exampleblock}
		
		\pp
		
		\begin{block}{6. Exponente cero}
			\[
			a^0 = 1, \qquad a\neq 0
			\]
		\end{block}
		
		\begin{exampleblock}{Ejemplo}
			\[
			7^0 = 1
			\]
		\end{exampleblock}
	\end{frame}
	
	% ----------------------------------------------------------
	\begin{frame}{Propiedades de las potencias IV}
		\begin{block}{7. Exponente uno}
			\[
			a^1 = a
			\]
		\end{block}
		
		\begin{exampleblock}{Ejemplo}
			\[
			9^1 = 9
			\]
		\end{exampleblock}
		
		\pp
		
		\begin{block}{8. Exponente negativo}
			\[
			a^{-n}=\frac{1}{a^n}, \qquad a\neq 0
			\]
		\end{block}
		
		\begin{exampleblock}{Ejemplo}
			\[
			2^{-3}=\frac{1}{2^3}=\frac{1}{8}
			\]
		\end{exampleblock}
	\end{frame}
	
	% ----------------------------------------------------------
	\begin{frame}{Propiedades de las potencias V}
		\begin{block}{9. Base igual a uno}
			\[
			1^n=1
			\]
		\end{block}
		
		\begin{exampleblock}{Ejemplo}
			\[
			1^{15}=1
			\]
		\end{exampleblock}
		
		\pp
		
		\begin{block}{10. Base igual a cero}
			\[
			0^n=0, \qquad n>0
			\]
		\end{block}
		
		\begin{exampleblock}{Ejemplo}
			\[
			0^4=0
			\]
		\end{exampleblock}
	\end{frame}
	
	% ----------------------------------------------------------
	\begin{frame}{Errores frecuentes}
		\begin{alertblock}{Error 1}
			\[
			a^m + a^n \neq a^{m+n}
			\]
		\end{alertblock}
		
		\begin{exampleblock}{Ejemplo}
			\[
			2^2+2^3 = 4+8=12
			\]
			pero
			\[
			2^{2+3}=2^5=32
			\]
		\end{exampleblock}
		
		\pp
		
		\begin{alertblock}{Error 2}
			\[
			(a+b)^n \neq a^n+b^n
			\]
		\end{alertblock}
		
		\begin{exampleblock}{Ejemplo}
			\[
			(2+3)^2=5^2=25
			\]
			pero
			\[
			2^2+3^2=4+9=13
			\]
		\end{exampleblock}
	\end{frame}
	
	% ----------------------------------------------------------
	\begin{frame}{Patrones en las potencias de 2}
		Observe la secuencia:
		\[
		2^1=2,\quad 2^2=4,\quad 2^3=8,\quad 2^4=16,\quad 2^5=32
		\]
		
		\begin{block}{Observación}
			Cada vez que el exponente aumenta en 1, el valor se multiplica por 2.
		\end{block}
		
		\pp
		
		\begin{ideaBox}
			Las potencias permiten describir \destacar{crecimientos rápidos} y \destacar{regularidades numéricas}.
		\end{ideaBox}
	\end{frame}
	
	% ----------------------------------------------------------
	\begin{frame}{Patrones y relaciones en las potencias de 10}
		\begin{columns}[T,totalwidth=\textwidth]
			\begin{column}{0.48\textwidth}
				\begin{block}{Exponente positivo}
					\[
					10^1 = 10
					\]
					\[
					10^2 = 100
					\]
					\[
					10^3 = 1000
					\]
					\[
					10^4 = 10000
					\]
				\end{block}
			\end{column}
			
			\begin{column}{0.48\textwidth}
				\begin{exampleblock}{Exponente negativo}
					\[
					10^0 = 1
					\]
					\[
					10^{-1} = 0{,}1
					\]
					\[
					10^{-2} = 0{,}01
					\]
					\[
					10^{-3} = 0{,}001
					\]
				\end{exampleblock}
			\end{column}
		\end{columns}
		
		\vspace{0.2cm}
		\begin{ideaBox}
			\textbf{Idea clave:} si el exponente sube, el número crece; si el exponente baja, el número disminuye.
		\end{ideaBox}
	\end{frame}
	
	% ----------------------------------------------------------
	\begin{frame}{Aplicación agronómica 1: semillas por parcela}
		\begin{block}{Situación}
			En una parcela experimental se siembran:
			\[
			2^5
			\]
			semillas por metro cuadrado.
			
			Si el terreno tiene:
			\[
			2^3
			\]
			metros cuadrados, determine el total de semillas.
		\end{block}
		
		\pp
		
		\begin{exampleblock}{Resolución}
			\[
			2^5\cdot2^3 = 2^{5+3}=2^8=256
			\]
			Se utilizan \textbf{256 semillas}.
		\end{exampleblock}
	\end{frame}
	
	% ----------------------------------------------------------
	\begin{frame}{Aplicación agronómica 2: crecimiento bacteriano}
		\begin{block}{Situación}
			Una colonia bacteriana del suelo se duplica cada hora.
			
			Si inicialmente hay:
			\[
			2^3
			\]
			bacterias, después de 4 horas habrá:
			\[
			2^3\cdot2^4
			\]
		\end{block}
		
		\pp
		
		\begin{exampleblock}{Resolución}
			\[
			2^3\cdot2^4=2^{7}=128
			\]
			Después de 4 horas hay \textbf{128 bacterias}.
		\end{exampleblock}
	\end{frame}
	
	% ----------------------------------------------------------
	\begin{frame}{Aplicación agronómica 3: fertilización}
		\begin{block}{Situación}
			Un fertilizante se aplica en dosis de:
			\[
			5\times10^2
			\]
			gramos por hectárea.
			
			Si el ensayo ocupa:
			\[
			10^3
			\]
			hectáreas experimentales, calcule la cantidad total.
		\end{block}
		
		\pp
		
		\begin{exampleblock}{Resolución}
			\[
			(5\times10^2)(10^3)=5\times10^5
			\]
			\[
			5\times10^5=500000
			\]
			Se requieren \textbf{500000 gramos}.
		\end{exampleblock}
	\end{frame}
	
	% ----------------------------------------------------------
	\begin{frame}{Aplicación agronómica 4: concentración de micronutrientes}
		\begin{block}{Situación}
			Una solución nutritiva contiene:
			\[
			3\times10^{-2}\ \text{g}
			\]
			de un micronutriente por litro.
			
			Si se preparan:
			\[
			10^3
			\]
			litros, ¿cuántos gramos del micronutriente se necesitan?
		\end{block}
		
		\pp
		
		\begin{exampleblock}{Resolución}
			\[
			(3\times10^{-2})(10^3)=3\times10^{1}
			\]
			\[
			3\times10^1=30
			\]
			Se necesitan \textbf{30 gramos}.
		\end{exampleblock}
	\end{frame}
	
	% ----------------------------------------------------------
	\begin{frame}{Ejercicios propuestos}
		Resuelva:
		\begin{enumerate}
			\item \(2^4\cdot2^3\)
			\item \((3^2)^3\)
			\item \(\dfrac{10^6}{10^2}\)
			\item \(\left(\dfrac{2}{5}\right)^2\)
			\item \(4^0\)
			\item \(3^{-2}\)
			\item \((2\cdot3)^2\)
			\item En un cultivo hay \(3^4\) plantas por sector y \(3^2\) sectores. ¿Cuántas plantas hay en total?
		\end{enumerate}
	\end{frame}
	
	% ----------------------------------------------------------
	\begin{frame}{Ejercicios con respuestas}
		\begin{multicols}{2}
			\begin{enumerate}
				\item \(2^4\cdot2^3=2^7=128\)
				\item \((3^2)^3=3^6=729\)
				\item \(\dfrac{10^6}{10^2}=10^4=10000\)
				\item \(\left(\dfrac{2}{5}\right)^2=\dfrac{4}{25}\)
				\item \(4^0=1\)
				\item \(3^{-2}=\dfrac{1}{9}\)
				\item \((2\cdot3)^2=2^2\cdot3^2=36\)
				\item \(3^4\cdot3^2=3^6=729\)
			\end{enumerate}
		\end{multicols}
	\end{frame}
	
	% ----------------------------------------------------------
	\begin{frame}{Síntesis}
		\begin{block}{Ideas clave}
			\begin{itemize}
				\item Las potencias representan multiplicaciones repetidas.
				\item La base indica el número que se repite.
				\item El exponente indica cuántas veces se repite.
				\item Las propiedades permiten simplificar cálculos de forma eficiente.
				\item En Agronomía, las potencias ayudan a modelar crecimiento, concentración y escalas.
			\end{itemize}
		\end{block}
		
		\pp
		
		\begin{ideaBox}
			\centering
			\textbf{Aprender propiedades de potencias permite calcular más rápido, cometer menos errores y comprender mejor fenómenos reales.}
		\end{ideaBox}
	\end{frame}
	
\end{document}